\documentclass[11pt,letterpaper]{article}

\usepackage[margin=0.85in]{geometry}
\usepackage{fontspec}
\usepackage{xcolor}
\usepackage{hyperref}
\usepackage{enumitem}
\usepackage{array}
\usepackage{longtable}
\usepackage{booktabs}
\usepackage{fancyhdr}
\usepackage{titlesec}
\usepackage{lastpage}
\usepackage{microtype}

\setmainfont{Times New Roman}
\setsansfont{Arial}
\hypersetup{
  colorlinks=true,
  linkcolor=blue!45!black,
  urlcolor=blue!45!black,
  pdftitle={Week 3 Comprehensive Lecture Notes},
  pdfauthor={EDCI 59100-016}
}

\pagestyle{fancy}
\fancyhf{}
\lhead{EDCI 59100-016}
\chead{Week 3}
\rhead{Finding, Reading, and Evaluating Literature}
\cfoot{\thepage}
\setlength{\headheight}{14pt}

\titleformat{\section}{\Large\bfseries\sffamily\color{blue!35!black}}{\thesection}{0.6em}{}
\titleformat{\subsection}{\large\bfseries\sffamily\color{blue!25!black}}{\thesubsection}{0.6em}{}
\titleformat{\subsubsection}{\normalsize\bfseries\sffamily}{\thesubsubsection}{0.6em}{}

\setlist[itemize]{topsep=3pt,itemsep=2pt,parsep=1pt,leftmargin=1.4em}
\setlist[enumerate]{topsep=3pt,itemsep=2pt,parsep=1pt,leftmargin=1.6em}
\renewcommand{\arraystretch}{1.18}
\setlength{\tabcolsep}{3pt}

\newcommand{\callout}[2]{
  \vspace{0.55em}
  \noindent\fbox{\parbox{0.965\linewidth}{\textbf{#1.} #2}}
  \vspace{0.55em}
}

\newcommand{\figurebox}[2]{
  \vspace{0.65em}
  \noindent\fbox{\parbox{0.965\linewidth}{\textbf{Figure placeholder: #1}\\#2}}
  \vspace{0.65em}
}

\newcommand{\searchstring}[1]{
  \begin{quote}
  \small\ttfamily\raggedright #1\par
  \end{quote}
}

\begin{document}

\begin{titlepage}
  \centering
  \vspace*{0.9in}
  {\Large EDCI 59100-016: Educational Research and Literature Review\par}
  \vspace{0.35in}
  {\Huge\bfseries Week 3 Comprehensive Lecture Notes\par}
  \vspace{0.2in}
  {\LARGE Finding, Reading, and Evaluating Literature\par}
  \vspace{0.55in}
  {\large Summer 2026\par}
  \vfill
  \begin{minipage}{0.82\linewidth}
  \centering
  \textbf{Major themes covered:} database searching, Boolean search strings, review-method alignment, search documentation, source screening, source evaluation, skimming scholarly articles, reference management software, AI-supported research workflows, discussion preparation, and the Search Results and Source Evaluation Matrix.
  \end{minipage}
  \vfill
  {\large Prepared from the saved Week 3 course materials in the local Brightspace course-content folder.\par}
\end{titlepage}

\tableofcontents
\newpage

\section{Learning Objectives}

After completing Week 3, students should be able to:

\begin{enumerate}
  \item Implement a search strategy using discipline-appropriate databases and structured search strings based on the Week 2 literature review plan.
  \item Translate research questions into searchable concept groups, keywords, synonyms, Boolean operators, quotation marks, and database-specific syntax.
  \item Document the search process transparently in a search log, with the amount of documentation matched to the selected review methodology.
  \item Screen candidate articles and make initial Keep, Maybe, or Exclude decisions using explicit relevance, quality, and method-fit criteria.
  \item Evaluate sources using criteria such as currency, relevance, authority, accuracy, purpose, methodology, evidence quality, bias, and contribution to the literature.
  \item Use skimming and scanning strategies to decide which scholarly articles deserve deeper reading.
  \item Organize sources in reference management software so that later synthesis and writing are more efficient.
  \item Explain how AI tools can support, but not replace, transparent searching, reading, screening, and documentation.
\end{enumerate}

\section{Topic Overview}

Week 3 marks a shift from designing a literature review plan to executing the early stages of the review. In Week 2, the emphasis was on clarifying a topic, selecting a review type, and planning a search. In Week 3, the emphasis moves to action: running searches, testing search terms, comparing databases, documenting decisions, screening candidate sources, and beginning to evaluate which studies are genuinely useful.

This topic matters because a literature review is only as credible as the process that produced it. A strong paper does not simply contain many sources. It shows that the researcher searched deliberately, selected sources for defensible reasons, evaluated evidence critically, and aligned every decision with the review methodology.

\callout{Key takeaway}{Searching is not a one-time task. It is an iterative research process. You test terms, examine results, revise the search, document what changed, and keep track of why articles were kept, marked as possible, or excluded.}

\subsection{How Week 3 Connects to Previous Work}

The Week 2 literature review plan provides the starting point:

\begin{itemize}
  \item The research question defines the main concepts to search.
  \item The review type defines how broad, selective, or systematic the search should be.
  \item The initial search plan suggests databases, keywords, and possible criteria.
  \item The purpose statement clarifies what kinds of evidence will matter.
\end{itemize}

Week 3 then tests whether the plan actually works in databases. If a search returns thousands of unrelated records, the plan needs refinement. If a search returns almost nothing, the plan may be too narrow or may use terminology that the literature does not use. If the first relevant studies use different keywords than expected, those keywords should be added to the search log and tested.

\subsection{Real-World Relevance}

Transparent literature searching is essential in graduate research, academic publishing, grant writing, curriculum design, health research, engineering reviews, policy work, and industry research. A project team may need to know what has already been tried, what evidence is strong, where results conflict, and where gaps remain. These decisions depend on the same skills practiced this week: searching, documenting, screening, evaluating, and organizing sources.

\figurebox{Literature review workflow}{A process diagram should show the flow from research question to concept groups, databases, search strings, search log, screening, source evaluation matrix, preliminary included studies, and later synthesis. The diagram should also show feedback loops from screening back to revised keywords and search strings.}

\section{The Week 3 Research Workflow}

The recommended weekly flow is:

\begin{enumerate}
  \item Search within databases for the topic.
  \item Read and evaluate candidate literature.
  \item Explore reference management software and citation managers.
  \item Apply algorithmic thinking when using AI tools for literature review work.
  \item Complete the discussion on screening decisions and search challenges.
  \item Complete the Search Results and Source Evaluation Matrix.
  \item Submit the final literature review plan.
\end{enumerate}

These tasks are connected. The search log records what was done. The screening table records what was found and what decision was made. The source evaluation matrix records what useful information each source contributes. The discussion gives a place to explain search challenges, compare strategies, and reflect on how method affects decisions.

\section{Searching Within Databases}

\subsection{Purpose of the Search Stage}

The purpose of Week 3 searching is to apply the plan to actual databases. At this stage, students are not expected to have the final perfect corpus of sources. They are expected to test and improve the process.

Before beginning, gather:

\begin{itemize}
  \item the research question or purpose statement;
  \item the list of keywords and synonyms from the literature review plan;
  \item the selected review type;
  \item the initial database plan;
  \item any inclusion and exclusion criteria already drafted.
\end{itemize}

\callout{Important}{The search process should be treated as evidence. A search log is not busywork. It helps show how the literature was located, how the strategy changed, and whether the review is methodologically credible.}

\subsection{Choosing Databases}

Different topics belong in different scholarly ecosystems. A database is not just a search box; it is a collection with its own disciplinary coverage, indexing rules, controlled vocabulary, filters, and search syntax.

\begin{longtable}{@{}p{0.21\linewidth}p{0.26\linewidth}p{0.47\linewidth}@{}}
\toprule
\textbf{Topic area} & \textbf{Useful databases} & \textbf{Why they fit}\\
\midrule
Education & ERIC, PsycINFO & ERIC focuses on education research; PsycINFO captures psychology, learning, cognition, motivation, behavior, and educational psychology.\\
Technology and computing & IEEE Xplore, ACM Digital Library & These databases capture computer science, engineering, systems, artificial intelligence, human-computer interaction, and technical conference literature.\\
Health and medicine & PubMed & PubMed is central for biomedical and health sciences literature and supports structured medical searching.\\
Broad interdisciplinary topics & Scopus, Web of Science, Google Scholar, OneSearch & These tools capture research across fields and can be useful when a topic crosses disciplinary boundaries.\\
\bottomrule
\end{longtable}

To decide where to search, ask:

\begin{itemize}
  \item Where is research in this field typically published?
  \item Does the database include the journals, conferences, or disciplines that matter for the topic?
  \item Do the articles already found cite sources from certain journals or fields repeatedly?
  \item Do Purdue library subject guides recommend discipline-specific databases?
  \item If uncertain, what happens when one discipline-specific database is compared with Google Scholar?
\end{itemize}

\subsection{Database Searching Is Testing, Not Finalizing}

The first database choice is not always the final choice. A database may be too broad, too narrow, too technical, too clinical, or too weak for a topic. Students should test databases by examining the first page or two of results and asking whether the records consistently align with the research question.

For example, a search about AI feedback in college writing may produce useful education studies in ERIC, psychology-oriented studies in PsycINFO, technical writing studies in Scopus, and computer-science model papers in ACM or IEEE. The correct choice depends on whether the review is about teaching practice, learner outcomes, writing assessment, model design, or the broad state of the field.

\section{Building Search Strings}

\subsection{Why Search Strings Matter}

Academic databases do not understand a full research question the way a human reader does. They match words, phrases, metadata, subject terms, abstracts, titles, authors, and indexed fields. A search string translates the research question into a form the database can process.

If a search is too simple, it can return too many unrelated results. If a search is too narrow, it can miss important studies. Search strings help control the search by combining the major concepts of the research question.

\subsection{Step-by-Step Search String Construction}

\begin{enumerate}
  \item Identify two or three key concepts from the research question.
  \item List synonyms, abbreviations, and related terms for each concept.
  \item Combine similar terms with \texttt{OR}.
  \item Combine different concept groups with \texttt{AND}.
  \item Use quotation marks for exact phrases.
  \item Use parentheses to group terms.
  \item Test the string in a database and inspect the first 5 to 10 results.
  \item Revise one element at a time so you can tell what improved or weakened the search.
\end{enumerate}

\subsection{Example Search Strings}

For a topic about AI feedback in writing, possible strings include:

\searchstring{("artificial intelligence" OR AI) AND (feedback OR "automated feedback") AND writing}

\searchstring{AI AND writing AND ("learning outcomes" OR performance)}

The first string is broader because it includes multiple feedback terms. The second is narrower because it adds outcomes or performance. Neither string is automatically correct. Each must be tested against the review purpose.

\subsection{How Databases Read the Search}

Most academic databases treat search syntax as formal instructions:

\begin{itemize}
  \item \texttt{AND} requires both concept groups to appear.
  \item \texttt{OR} accepts any term within a group of similar terms.
  \item Quotation marks keep words together as a phrase.
  \item Parentheses control how terms are grouped.
  \item Field filters such as title, abstract, subject, or full text can change results dramatically.
  \item Small syntax changes can create large changes in the number and quality of records.
\end{itemize}

\callout{Common mistake}{Do not type the entire research question into a database and expect precise scholarly results. Break the question into searchable concepts and controlled combinations.}

\figurebox{Boolean search logic}{A diagram should show three concept circles: AI terms, writing terms, and feedback terms. The overlap where all three circles intersect represents the effect of AND. Smaller circles inside each concept group represent synonyms joined by OR.}

\section{Inclusion and Exclusion Criteria}

\subsection{Purpose}

Inclusion and exclusion criteria define what counts as a relevant source before the researcher becomes too attached to individual articles. Criteria help make screening consistent and defensible.

Initial criteria may include:

\begin{itemize}
  \item source type, such as empirical study, theoretical article, review article, dissertation, conference paper, or grey literature;
  \item population, such as K-12 students, higher education students, teachers, patients, workers, or organizations;
  \item context, such as online learning, construction sites, healthcare settings, engineering classrooms, or workplace training;
  \item time frame, such as the last 5 years or 2015 to 2025;
  \item method, such as qualitative, quantitative, mixed-methods, design-based research, experimental, or conceptual;
  \item language, publication status, peer-review status, and geographic scope;
  \item connection to the chosen review type.
\end{itemize}

\subsection{Keep, Maybe, Exclude}

Early screening often uses three practical categories:

\begin{description}[leftmargin=1.15in,style=nextline]
  \item[Keep] The source clearly matches the topic, criteria, and review purpose.
  \item[Maybe] The source is related but has an uncertain fit. It might differ by population, method, field, time period, or focus.
  \item[Exclude] The source does not fit the topic, criteria, review type, or quality expectations.
\end{description}

A Maybe source is not a failure. It is a normal part of early searching. In fact, Maybe sources are useful because they reveal boundary problems. If many Maybe sources are about a different population, the search may need more specific population terms. If many Maybe sources are technical model papers but the review is about learning outcomes, the search may need education-specific terms.

\section{Method-Specific Search Expectations}

The selected review methodology changes how the search should be planned, documented, and judged. The same article can be a strong Keep for one review type and a Maybe or Exclude for another.

\subsection{Narrative Review}

A narrative review explains a topic by drawing together important, relevant, and often influential literature. It does not aim to retrieve every possible study.

\textbf{Search emphasis:}

\begin{itemize}
  \item Begin broad, then refine toward studies that explain the topic clearly.
  \item Include foundational works and recent research.
  \item Use citation chaining to identify influential works.
  \item Prioritize sources that help tell the scholarly story of the topic.
\end{itemize}

\textbf{Documentation expectation:} Narrative reviews do not require rigid exhaustive search protocols, but the writer should still explain how key studies were identified.

\textbf{Guiding question:} What studies help explain this topic clearly?

\subsection{Integrative Review}

An integrative review brings together multiple forms of evidence and perspectives. It may include quantitative studies, qualitative studies, theoretical or conceptual work, and different methodological traditions.

\textbf{Search emphasis:}

\begin{itemize}
  \item Use multiple search strings and databases.
  \item Intentionally capture variation in study types and perspectives.
  \item Avoid collecting many similar studies while missing conceptual or methodological diversity.
  \item Look for patterns across heterogeneous sources.
\end{itemize}

\textbf{Documentation expectation:} Document enough of the search process to show how variation across study types and perspectives was intentionally captured.

\textbf{Guiding question:} What range of evidence is needed to understand patterns across studies?

\subsection{Theoretical Review}

A theoretical review focuses on ideas, concepts, models, frameworks, and the development of theoretical relationships. It may include fewer sources than other review types, but each source should be influential and conceptually important.

\textbf{Search emphasis:}

\begin{itemize}
  \item Search for key concepts, theorists, models, frameworks, and foundational terms.
  \item Prioritize seminal works and high-impact theoretical papers.
  \item Do not limit the search only to empirical studies.
  \item Use phrase searching, Boolean operators, proximity searching, and document-type filters when available.
\end{itemize}

\textbf{Documentation expectation:} Explain how selected works represent the development of the theory or concept.

\textbf{Guiding question:} What works help explain how this idea, theory, or framework developed?

\subsection{Scoping Review}

A scoping review maps the range, nature, and extent of literature on a topic. It is especially useful when a field is emerging, interdisciplinary, conceptually varied, or not yet ready for a narrow systematic review.

\textbf{Search emphasis:}

\begin{itemize}
  \item Search broadly across databases, journals, conferences, and source types.
  \item Capture different populations, contexts, methods, and study types.
  \item Include grey literature when it is relevant to the review purpose.
  \item Use citation chaining to expand from seed articles.
  \item Consider search updates because literature can become outdated quickly.
\end{itemize}

\textbf{Documentation expectation:} PRISMA-ScR expects exact search strings, dates of searches, number of results per database, and a flow of records through identification, screening, and inclusion.

\textbf{Guiding question:} What has been studied, in what ways, and what has not been studied?

\subsection{Systematic Review}

A systematic review follows a predefined, transparent, and replicable process. It usually has narrow criteria and aims to answer a focused question.

\textbf{Search emphasis:}

\begin{itemize}
  \item Define search terms in advance.
  \item Search multiple databases.
  \item Use clear inclusion and exclusion criteria.
  \item Record how many records were found, removed, screened, excluded, and included.
  \item Consider grey literature to reduce publication bias.
  \item Translate search strings for each database rather than pasting the same syntax everywhere.
\end{itemize}

\textbf{Documentation expectation:} Full documentation is required: database names, interfaces, search dates, full search strings, modifications, limits, filters, searcher information, results, screening counts, full-text exclusions, and final included studies.

\textbf{Guiding question:} Can someone else follow this process step by step and reach a comparable result set?

\begin{longtable}{@{}p{0.15\linewidth}p{0.27\linewidth}p{0.30\linewidth}p{0.20\linewidth}@{}}
\toprule
\textbf{Review type} & \textbf{Search goal} & \textbf{Best evidence to prioritize} & \textbf{Documentation level}\\
\midrule
Narrative & Explain the topic through key literature & Foundational works, recent studies, influential perspectives & Moderate; explain strategy and key choices\\
Integrative & Compare patterns across varied evidence & Quantitative, qualitative, theoretical, and conceptual sources & Moderate to high; show intentional variation\\
Theoretical & Trace concepts, frameworks, and models & Seminal theoretical works and conceptual papers & Moderate; explain conceptual lineage\\
Scoping & Map the extent and range of literature & Broad sources across contexts, populations, methods, and grey literature & High; PRISMA-ScR style tracking\\
Systematic & Answer a focused question with replicable searching & Studies meeting predefined criteria & Very high; PRISMA-style tracking\\
\bottomrule
\end{longtable}

\section{Citation Chaining and Snowballing}

Citation chaining helps identify literature that may not appear through keyword searches alone.

\begin{description}[leftmargin=1.35in,style=nextline]
  \item[Backward snowballing] Looking at the references of a relevant seed article to find earlier foundational studies.
  \item[Forward snowballing] Looking up who has cited a seed article in Google Scholar, Scopus, or Web of Science to find newer work that built on it.
\end{description}

Snowballing is useful in all review types, but its role varies. A narrative or theoretical review may use it to find influential or foundational sources. A scoping or systematic review may use it as a supplementary search method that must be documented.

\callout{Exam tip}{Be able to explain the difference between backward and forward snowballing and why both can improve a literature search.}

\section{Practicing Search Refinement}

The Week 3 practice task asks students to compare and refine a search.

\begin{enumerate}
  \item Choose one search string.
  \item Run two searches: one in a university database and one in Google Scholar.
  \item Compare the first 5 to 10 results.
  \item Ask what types of sources appear.
  \item Ask whether the sources are similar or different.
  \item Decide which search is more useful for the topic.
  \item Make one change: adjust a keyword, add specificity, or refine the focus.
  \item Run the search again and judge whether it improved.
\end{enumerate}

\subsection{How to Judge Improvement}

Improvement does not always mean fewer results. It means better alignment between the search output and the review purpose.

Ask:

\begin{itemize}
  \item Do the results consistently relate to the topic?
  \item Do they reflect the selected review method?
  \item Do they begin to show patterns across studies?
  \item Are the same irrelevant topics appearing repeatedly?
  \item Are key expected concepts missing?
  \item Does the search capture both foundational and recent work when needed?
\end{itemize}

\subsection{Method-Specific Refinement Questions}

\begin{itemize}
  \item Narrative review: Are these studies key and influential? Do they help explain the topic clearly?
  \item Integrative review: Do the results include different types of studies, or are they too similar?
  \item Theoretical review: Are the results finding concepts, frameworks, and foundational ideas, or mostly empirical reports?
  \item Scoping review: Do the results reflect a range of contexts, populations, and study types?
  \item Systematic review: Can the inclusion or exclusion of each result be justified by predefined criteria?
\end{itemize}

\section{Documenting the Search}

\subsection{The Search Log}

A search log records the search process in a transparent way. It helps the researcher remember what was tried, avoid repeating ineffective searches, explain decisions, and demonstrate methodological rigor.

A strong search log includes:

\begin{itemize}
  \item search number;
  \item database or platform used;
  \item exact search string;
  \item filters or limits applied;
  \item number of results;
  \item notes on relevance, problems, and decisions.
\end{itemize}

\begin{longtable}{@{}p{0.06\linewidth}p{0.14\linewidth}p{0.30\linewidth}p{0.18\linewidth}p{0.10\linewidth}p{0.16\linewidth}@{}}
\toprule
\textbf{\#} & \textbf{Database} & \textbf{Search string} & \textbf{Filters} & \textbf{Results} & \textbf{Notes}\\
\midrule
1 & ERIC & \texttt{("artificial intelligence" OR AI) AND feedback AND writing} & Peer-reviewed; last 10 years & 1,245 & Too broad; included K-12 and higher education.\\
2 & ERIC & \texttt{("artificial intelligence" OR AI) AND feedback AND writing AND "higher education"} & Peer-reviewed; last 10 years & 312 & More focused; better population alignment.\\
3 & Google Scholar & Same core concepts adapted to platform & Relevance sorting or date filters as needed & Record count varies & Useful for citation tracing, but less reproducible.\\
\bottomrule
\end{longtable}

\subsection{Search Log Quality}

A weak search log only lists what was searched. A strong search log explains what the search taught the researcher. For example, the note ``too broad'' is useful, but it becomes stronger when it identifies why the search was too broad: wrong population, wrong field, too many technical sources, too many opinion pieces, too many unrelated uses of a term, or missing method language.

\callout{Important}{Always record exact strings before editing them. Once a useful search is changed, it can be difficult to reconstruct what produced the earlier result set.}

\section{Screening Candidate Sources}

\subsection{Initial Screening Notes}

Screening asks whether each source belongs in the developing review. At the early stage, screening does not require a full detailed reading of every article. It uses titles, abstracts, keywords, headings, introductions, methods, conclusions, and sometimes reference lists to decide whether the source deserves further attention.

A screening table should include:

\begin{itemize}
  \item article title;
  \item Keep, Maybe, or Exclude decision;
  \item reason for the decision;
  \item explanation of method fit.
\end{itemize}

\begin{longtable}{@{}p{0.26\linewidth}p{0.12\linewidth}p{0.29\linewidth}p{0.27\linewidth}@{}}
\toprule
\textbf{Article title} & \textbf{Decision} & \textbf{Why} & \textbf{Fits method?}\\
\midrule
AI Feedback in College Writing Courses (2022) & Keep & Direct match to topic; higher education focus; empirical findings. & Integrative review: provides empirical evidence that can be compared with other study types.\\
AI Tools for K-12 Writing Support (2019) & Maybe & Relevant topic but different population. & Integrative review: may add variation across contexts if the review includes K-12 comparison.\\
General AI Productivity Tools in Business (2021) & Exclude & AI topic but not education, writing, feedback, or target context. & Does not fit the review purpose.\\
\bottomrule
\end{longtable}

\subsection{How Method Changes Screening}

The same source can be evaluated differently depending on the review type.

\begin{longtable}{@{}p{0.17\linewidth}p{0.17\linewidth}p{0.60\linewidth}@{}}
\toprule
\textbf{Method} & \textbf{Likely decision} & \textbf{Reasoning}\\
\midrule
Narrative & Keep & The article explains key ideas clearly and is widely cited or highly relevant.\\
Integrative & Keep & The article adds empirical evidence that can be compared with qualitative, theoretical, or conceptual work.\\
Theoretical & Maybe & The article is useful only if it connects to a theory, model, framework, or conceptual debate.\\
Scoping & Keep & The article represents one context, population, method, or dimension within the broader literature landscape.\\
Systematic & Keep or Exclude & The decision depends strictly on predefined criteria such as population, design, date range, intervention, outcome, and publication type.\\
\bottomrule
\end{longtable}

\section{Evaluating Sources}

\subsection{When Evaluation Happens}

Source evaluation happens at several stages:

\begin{enumerate}
  \item when the source first appears in search results;
  \item when the abstract and introduction are skimmed;
  \item when the source is read more carefully;
  \item when it is compared with other sources during synthesis;
  \item when it is finally used in the literature review.
\end{enumerate}

The goal is to ensure the review uses literature that is relevant, accurate, credible, and appropriately interpreted.

\subsection{Big 5 Criteria}

The Big 5 criteria are:

\begin{description}[leftmargin=1.35in,style=nextline]
  \item[Currency] When was the source published? Is it current enough for the field?
  \item[Coverage/Relevance] How closely does it connect to the research question?
  \item[Authority] Who wrote it? Are the authors credible for this topic?
  \item[Accuracy] Are claims supported by evidence? Are there obvious errors?
  \item[Objectivity/Purpose] Does the source have a bias, agenda, or purpose that affects interpretation?
\end{description}

For medical, scientific, and technology topics, sources from the last 5 years are often preferred because knowledge changes quickly. For less time-sensitive humanities or historical topics, sources from the last 5 to 10 years may be acceptable, and older foundational sources may remain essential.

\subsection{CRAAP Test}

The CRAAP Test is a mnemonic that overlaps with the Big 5:

\begin{description}[leftmargin=1.15in,style=nextline]
  \item[Currency] Is the source current enough?
  \item[Relevance] Does it answer or inform the research question?
  \item[Authority] Is the author or organization credible?
  \item[Accuracy] Are the claims evidence-based and verifiable?
  \item[Purpose] Why was the source created? To inform, persuade, sell, entertain, or advocate?
\end{description}

\subsection{Initial Evaluation Questions}

During the first pass, ask:

\begin{itemize}
  \item When was the source published?
  \item Is it relevant to the topic?
  \item What are the author qualifications?
  \item Is the source scholarly or peer-reviewed?
  \item Are sources cited to support claims?
  \item Does the journal, website, or organization have a clear bias?
  \item Does the summary or introduction use emotional or biased language?
  \item Are there obvious spelling, grammar, or factual problems?
\end{itemize}

\subsection{Deeper Evaluation During Reading}

Once a source passes initial evaluation, read it with more precise questions:

\begin{enumerate}
  \item Is there visible bias in the work?
  \item How was the research conducted?
  \item What are the strengths and weaknesses of the methodology?
  \item How does the author justify the conclusions?
  \item What similarities does this source share with other articles?
  \item How does this source differ from other articles?
  \item What new knowledge, evidence, method, or insight does it contribute?
  \item What unanswered questions or future research opportunities does it leave?
\end{enumerate}

\callout{Common mistake}{A peer-reviewed article is not automatically strong for every literature review. It still needs to fit the question, method, criteria, and purpose of the review.}

\section{Reading Scholarly Articles Efficiently}

\subsection{Why Strategic Reading Matters}

Literature review work requires reading many sources under time constraints. Reading every article from the first word to the last word is often inefficient during early screening. Strategic reading helps decide which sources deserve deep attention.

\subsection{Recommended Reading Sequence}

The saved reading strategy materials recommend skipping around at first:

\begin{enumerate}
  \item Read the title and abstract.
  \item Skim the introduction.
  \item Jump to the conclusion or discussion.
  \item Look at results, figures, tables, charts, graphs, or images.
  \item If the source is relevant, then read it from start to finish.
  \item Take notes, summarize sections, and record citation information.
  \item Check the reference list for additional related sources.
\end{enumerate}

This sequence works because the abstract, introduction, discussion, and conclusion usually explain the article's purpose, context, findings, and contribution. The methods and results are critical for careful evaluation, but they may not be the best starting point during early relevance screening.

\subsection{The Outside-In Article Strategy}

Scholarly articles often put the big-picture information on the outside:

\begin{itemize}
  \item abstract;
  \item introduction;
  \item discussion;
  \item conclusion.
\end{itemize}

The technical core is often inside:

\begin{itemize}
  \item methodology;
  \item instruments;
  \item data collection;
  \item statistical analysis;
  \item results details.
\end{itemize}

The outside sections help decide relevance. The inside sections help evaluate validity, replicability, and strength of evidence.

\section{Skimming and Scanning}

\subsection{Skimming}

Skimming means reading quickly to understand the main ideas, structure, and argument without focusing on every word. It helps researchers decide whether a source is relevant enough for deeper reading.

Types of skimming include:

\begin{description}[leftmargin=1.55in,style=nextline]
  \item[Preview skimming] Looking at titles, headings, abstracts, and introductions before deciding whether to read closely.
  \item[Overview skimming] Moving through the full article quickly to understand its structure, flow, and argument.
  \item[Selective skimming] Looking for specific terms, sections, or concepts related to the research question.
  \item[Comprehensive skimming] Combining speed with light analysis by noticing evidence, examples, and data while avoiding full deep reading.
\end{description}

\subsection{Scanning}

Scanning means looking for specific information such as dates, names, measures, populations, theories, keywords, or statistics. Skimming gives the big picture. Scanning locates precise details.

\subsection{A Practical Skimming Procedure}

\begin{enumerate}
  \item Start with the title and abstract. Ask whether the source aligns with the research question.
  \item Read the introduction and conclusion to understand the purpose and outcome.
  \item Review headings and subheadings to understand organization.
  \item Examine visuals and data displays.
  \item Read the first and last sentences of important paragraphs.
  \item Mark keywords, definitions, theories, and method terms.
  \item Review references for possible citation chaining.
\end{enumerate}

\callout{Key takeaway}{Skimming is not lazy reading. It is a screening strategy. Deep reading should be reserved for sources that pass initial relevance and quality checks.}

\section{Reference Management Software}

\subsection{Why Use a Citation Manager}

Reference management software supports the literature review workflow by storing citations, PDFs, notes, tags, folders, metadata, and citation exports. It reduces the risk of losing sources and helps organize evidence before writing.

Common tools include:

\begin{itemize}
  \item Zotero;
  \item Mendeley;
  \item EndNote;
  \item RefWorks;
  \item database export files such as RIS and BibTeX.
\end{itemize}

\subsection{Recommended Organization}

For Week 3, a practical source library can include:

\begin{itemize}
  \item a folder for all raw search exports;
  \item a folder for Keep sources;
  \item a folder for Maybe sources;
  \item a folder for Excluded-after-screening sources if the review requires tracking;
  \item tags for method, population, context, theory, database, and screening decision;
  \item notes that summarize why each source matters.
\end{itemize}

\subsection{Connection to Later Writing}

The goal of reference management is not just citation formatting. The deeper purpose is retrieval. When writing the literature review, the researcher needs to quickly find studies by theme, method, finding, population, or limitation. A well-organized library makes synthesis easier.

\section{AI Tools and Algorithmic Thinking}

\subsection{What AI Can Support}

AI tools can help with parts of the literature review workflow:

\begin{itemize}
  \item brainstorming synonyms and concept groups;
  \item translating a research question into possible search strings;
  \item summarizing long abstracts or articles;
  \item identifying recurring concepts across notes;
  \item checking whether a screening rationale is clear;
  \item generating draft tables for organization;
  \item creating plain-language explanations of difficult methods.
\end{itemize}

\subsection{What AI Cannot Replace}

AI tools cannot replace the researcher's responsibility to:

\begin{itemize}
  \item verify every source against the original article;
  \item record the actual database search string and result count;
  \item decide inclusion and exclusion based on the selected review method;
  \item evaluate methodology and evidence quality;
  \item avoid fabricated citations or unsupported claims;
  \item preserve transparency and replicability.
\end{itemize}

\callout{Important}{AI output should be treated as a draft aid, not as evidence. The evidence is the database record, article, method, results, and documented screening decision.}

\section{Assignment Focus: Search Results and Source Evaluation Matrix}

The Week 3 assignment asks students to submit a working Searching and Evaluating Literature document. The assignment focuses on documenting and reflecting on the process for identifying and evaluating literature.

\subsection{What the Assignment Should Demonstrate}

The assignment should show that the student is:

\begin{itemize}
  \item developing an effective search strategy;
  \item applying inclusion and exclusion criteria;
  \item making thoughtful screening decisions;
  \item beginning to identify patterns across sources;
  \item revising the process based on feedback and search results.
\end{itemize}

\subsection{Required Components}

\begin{enumerate}
  \item \textbf{Search Log:} databases used, search strings, filters, number of results, and notes on search effectiveness.
  \item \textbf{Screening Decisions:} sources reviewed, Keep/Maybe/Exclude decisions, and brief explanations tied to criteria.
  \item \textbf{Preliminary Set of Included Studies:} sources likely to remain in the review, even if the final set changes.
  \item \textbf{Source Evaluation Matrix:} article notes that include main idea, method, key findings, and why the source matters.
\end{enumerate}

The expected scope is at least 10 sources across screening and evaluation. These sources may include Keep, Maybe, and Exclude decisions. The purpose is to show process and reasoning, not merely to collect a high number of citations.

\subsection{Evaluation Criteria}

The assignment is evaluated on:

\begin{itemize}
  \item quality of the search process;
  \item clarity of screening decisions and reasoning;
  \item development of source evaluation notes;
  \item use of criteria and method alignment;
  \item overall clarity and organization.
\end{itemize}

\begin{longtable}{@{}p{0.20\linewidth}p{0.20\linewidth}p{0.22\linewidth}p{0.23\linewidth}p{0.07\linewidth}@{}}
\toprule
\textbf{Criterion} & \textbf{Excellent} & \textbf{Proficient} & \textbf{Developing} & \textbf{Points}\\
\midrule
Search log & Multiple searches documented with refinement and thoughtful notes & Searches documented with some explanation & Limited searches or minimal explanation & 5\\
Screening decisions & Clear Keep/Maybe/Exclude decisions using criteria & Decisions provided with some reasoning & Decisions unclear or inconsistent & 5\\
Source set and evaluation & Robust set of about 10 or more sources with developed notes & Adequate set and some detail & Fewer sources or limited evaluation & 5\\
Criteria and method alignment & Decisions clearly connect to criteria and review method & Some connection present & Connection limited or unclear & 3\\
Clarity and organization & Well organized and easy to follow & Generally clear & Somewhat difficult to follow & 2\\
\bottomrule
\end{longtable}

\section{Worked Example: Refining a Search}

Suppose the research question is:

\begin{quote}
How does AI-supported automated feedback affect writing development in higher education?
\end{quote}

\subsection{Step 1: Identify Concept Groups}

\begin{longtable}{@{}p{0.25\linewidth}p{0.70\linewidth}@{}}
\toprule
\textbf{Concept} & \textbf{Possible terms}\\
\midrule
AI & artificial intelligence, AI, machine learning, automated, generative AI\\
Feedback & feedback, automated feedback, formative feedback, writing feedback\\
Writing & writing, composition, academic writing, writing instruction\\
Context & higher education, college, university, undergraduate, graduate\\
\bottomrule
\end{longtable}

\subsection{Step 2: Build a Broad Search}

\searchstring{("artificial intelligence" OR AI) AND (feedback OR "automated feedback") AND writing}

This string may be too broad because it can retrieve K-12 studies, general tool discussions, or non-education writing contexts.

\subsection{Step 3: Add Context}

\searchstring{("artificial intelligence" OR AI) AND (feedback OR "automated feedback") AND writing AND ("higher education" OR college OR university)}

This version is more focused. If it removes too many important studies, the researcher can test whether some articles use ``postsecondary'' instead of ``higher education'' or whether a database subject heading captures college students more effectively.

\subsection{Step 4: Screen Sample Results}

The first 10 results should be inspected for source type, population, topic match, method, and review-method fit. The researcher should not silently choose only the articles that look easy to use. Every Keep, Maybe, or Exclude decision should be explainable.

\section{Source Evaluation Matrix}

A source evaluation matrix transforms articles into usable synthesis notes.

\begin{longtable}{@{}p{0.18\linewidth}p{0.20\linewidth}p{0.14\linewidth}p{0.25\linewidth}p{0.18\linewidth}@{}}
\toprule
\textbf{Article} & \textbf{Main idea} & \textbf{Method} & \textbf{Key findings} & \textbf{Why it matters}\\
\midrule
Article A & AI feedback system in college writing & Mixed methods & Students valued speed, but quality varied by task & Useful for feedback benefits and limitations\\
Article B & Automated writing evaluation in K-12 & Quantitative & Scores improved, but population differs from target & Maybe source for context comparison\\
Article C & Conceptual framework for AI writing support & Theoretical & Defines feedback timing, agency, and revision support & Useful for theoretical framing\\
\bottomrule
\end{longtable}

Good matrix notes are brief but specific. They should not copy the abstract. They should extract the information needed for the literature review: what the study did, what it found, how strong the evidence is, and how it connects to the review purpose.

\section{Common Misconceptions}

\begin{description}[leftmargin=1.55in,style=nextline]
  \item[More sources means better review.] Quality, fit, method alignment, and synthesis matter more than raw count.
  \item[Google Scholar is enough.] Google Scholar is useful, especially for citation chasing, but it is less controlled and less reproducible than discipline databases.
  \item[A perfect search exists.] Searching is iterative. A documented, reasonable, well-refined process is the goal.
  \item[Maybe sources are useless.] Maybe sources clarify boundaries and can reveal whether criteria need refinement.
  \item[AI summaries can replace reading.] AI summaries must be verified against the original article and cannot replace evaluation of methods and evidence.
  \item[All review types require the same search.] Narrative, integrative, theoretical, scoping, and systematic reviews have different search and documentation expectations.
\end{description}

\section{Important Notes for Studying}

\callout{Important}{Review methodology controls search behavior. Always ask what the review type requires before judging whether the search is strong.}

\callout{Exam tip}{Know the difference between AND and OR. AND narrows by requiring multiple concepts. OR broadens by allowing synonyms or related terms within the same concept group.}

\callout{Common mistake}{Do not apply filters without documenting them. Date limits, language limits, peer-reviewed filters, and subject filters can all change the evidence base.}

\callout{Key takeaway}{The Week 3 deliverable is a working document. Sources and decisions can change later, but the current process should be clear, organized, and defensible.}

\section{Summary}

Week 3 develops the practical mechanics of literature review work. Students move from a plan to a documented search process. They choose databases based on field and topic, build search strings from concept groups, test Boolean logic, compare database outputs, refine keywords, define inclusion and exclusion criteria, and begin screening sources.

The selected review method determines the search goal. Narrative reviews prioritize key explanatory works. Integrative reviews intentionally include varied evidence. Theoretical reviews trace concepts and frameworks. Scoping reviews map the range of literature. Systematic reviews require a predefined, transparent, and replicable process.

Source evaluation happens at multiple points. The Big 5 and CRAAP criteria help judge currency, relevance, authority, accuracy, and purpose. Deeper reading adds questions about bias, methodology, conclusions, similarities, differences, contributions, and research gaps.

Strategic reading is essential. Skimming, scanning, and outside-in article reading help researchers decide which sources deserve close attention. Reference management software supports organization, and AI tools can assist with brainstorming, summarizing, and structuring, but the researcher remains responsible for verification and documentation.

\section{Review Questions}

\subsection{Conceptual Questions}

\begin{enumerate}
  \item Why is a search log important for a literature review?
  \item How does a search string differ from a research question?
  \item Why does \texttt{OR} usually broaden a search while \texttt{AND} usually narrows it?
  \item What is the difference between initial screening and deeper source evaluation?
  \item How does a scoping review search differ from a narrative review search?
  \item Why might the same article be Keep for one review type and Maybe for another?
  \item What are the Big 5 criteria for source evaluation?
  \item How does the CRAAP Test help evaluate sources?
  \item Why should AI-generated summaries be checked against the original article?
  \item What is the difference between backward and forward snowballing?
\end{enumerate}

\subsection{Short-Answer Practice With Answers}

\textbf{Question 1.} A student searches \texttt{AI writing feedback} and gets thousands of unrelated results. What are two improvements?

\textbf{Answer.} The student can group synonyms with Boolean operators, such as:
\searchstring{("artificial intelligence" OR AI) AND (feedback OR "automated feedback") AND writing}
The student can also add context terms such as \texttt{"higher education"} or filter by peer-reviewed sources and date range if those limits match the review criteria.

\textbf{Question 2.} Why should a systematic review document database-specific syntax changes?

\textbf{Answer.} Different databases use different field codes, subject headings, truncation rules, and proximity operators. Documenting syntax changes makes the process transparent and helps another researcher understand how the same conceptual search was adapted across platforms.

\textbf{Question 3.} A source is relevant to the topic but uses a different population. Should it automatically be excluded?

\textbf{Answer.} No. It may be marked Maybe. The decision depends on the review method and criteria. An integrative or scoping review may use the source for variation or mapping, while a systematic review may exclude it if the population does not meet predefined criteria.

\subsection{Application Questions}

\begin{enumerate}
  \item Build a search string for a research question about virtual reality training and teacher preparation.
  \item Create three possible inclusion criteria and three exclusion criteria for a review about AI tutoring in higher education.
  \item Choose one article from your own search results and justify a Keep, Maybe, or Exclude decision.
  \item Compare how a narrative review and scoping review would treat the same broad topic.
  \item Create a one-row source evaluation matrix for one source you found this week.
\end{enumerate}

\section{Cheat Sheet}

\begin{longtable}{@{}p{0.25\linewidth}p{0.70\linewidth}@{}}
\toprule
\textbf{Concept} & \textbf{Rule of thumb}\\
\midrule
Database choice & Match the database to the discipline and topic.\\
Search string & Convert research questions into concept groups and synonyms.\\
\texttt{AND} & Narrows by requiring different concepts together.\\
\texttt{OR} & Broadens by including synonyms or related terms.\\
Quotation marks & Keep phrases together, such as \texttt{"automated feedback"}.\\
Parentheses & Group terms so the database reads the logic correctly.\\
Inclusion criteria & Define what belongs in the review.\\
Exclusion criteria & Define what does not belong in the review.\\
Keep & Clearly fits topic, criteria, and method.\\
Maybe & Related but uncertain fit; useful for boundary decisions.\\
Exclude & Does not fit topic, criteria, method, or quality expectations.\\
Big 5 & Currency, coverage/relevance, authority, accuracy, objectivity/purpose.\\
CRAAP & Currency, relevance, authority, accuracy, purpose.\\
Backward snowballing & Search the references of seed articles.\\
Forward snowballing & Search newer articles that cited seed articles.\\
Scoping review & Map the range of evidence and gaps.\\
Systematic review & Follow a predefined, transparent, replicable process.\\
Skimming & Read quickly for structure and main ideas.\\
Scanning & Search for specific information.\\
Reference manager & Store, tag, organize, annotate, and cite sources.\\
\bottomrule
\end{longtable}

\section{Source Materials Used}

These notes were prepared from the saved Week 3 course-content materials, including the local Week 3 overview, ``Searching Within Databases for Your Topic,'' ``Reading and Evaluating Literature,'' ``Tool Intro - Reference Management Software,'' ``AI Tools - AI Algorithmic Thinking for Literature Review,'' the discussion and assignment references visible in the Week 3 module, the local assignment prompt for the Search Results and Source Evaluation Matrix, and the saved external reading pages on evaluating sources and reading/skimming scholarly articles.

\end{document}
